% 2011 年湖北卷理科第 18 题：正三棱柱 ABC-A1B1C1，E 为 BC 中点，F 在 CC1 上。
% 来源 PDF：普通高考/2011/2011湖北理.pdf，SHA-256 ba721f03ae90f88887c5cde27ae98dbea165d33c73e5ba45b4e4f56b0638b075。
% 原页：第 2 页第 1 栏上部，书页码 668。示意。
\begin{tikzpicture}[>=stealth, semithick, font=\small, scale=0.95]
  \coordinate (A) at (0,0);
  \coordinate (B) at (3.0,0);
  \coordinate (C) at (1.5,1.0);
  \coordinate (A1) at (0,3.0);
  \coordinate (B1) at (3.0,3.0);
  \coordinate (C1) at (1.5,4.0);
  \coordinate (E) at ($(B)!0.5!(C)$);
  \coordinate (F) at ($(C)!0.45!(C1)$);
  \draw[thick] (A)--(B)--(C)--cycle;
  \draw[thick] (A1)--(B1)--(C1)--cycle;
  \draw[thick] (A)--(A1) (B)--(B1) (C)--(C1);
  \draw (A1)--(C);
  \draw (E)--(F);
  \draw (A)--(F);
  \draw (A)--(E);
  \fill (E) circle (1.2pt) node[below right] {$E$};
  \fill (F) circle (1.2pt) node[right] {$F$};
  \node[left] at (A) {$A$};
  \node[right] at (B) {$B$};
  \node[below] at (C) {$C$};
  \node[left] at (A1) {$A_1$};
  \node[right] at (B1) {$B_1$};
  \node[above] at (C1) {$C_1$};
\end{tikzpicture}
